% =====================================================================
%  DashyDashboard / Access Review — End-User Guide
%  Plain-language walkthrough of every role and what each button does.
%
%  Build:  pdflatex UserGuide.tex   (run twice for the table of contents)
%     or:  latexmk -pdf UserGuide.tex
%
%  Screenshots: drop PNG/JPG files into ./figures and replace the
%  \shotph{...} placeholders with \shot{filename}{caption}. Until then
%  the placeholders render as labelled empty boxes so the document is
%  complete and print-ready as-is.
% =====================================================================
\documentclass[11pt]{article}

\usepackage[a4paper,margin=2.4cm,headheight=14pt]{geometry}
\usepackage[T1]{fontenc}
\usepackage[utf8]{inputenc}
\usepackage{lmodern}
\usepackage{microtype}
\usepackage{xcolor}
\usepackage{graphicx}
\usepackage{enumitem}
\usepackage{booktabs}
\usepackage{array}
\usepackage{longtable}
\usepackage{tabularx}
\usepackage{titlesec}
\usepackage{fancyhdr}
\usepackage[most]{tcolorbox}
\usepackage{fontawesome5}
\usepackage[hidelinks]{hyperref}

% ---------- Brand palette ----------
\definecolor{brand}{HTML}{1F6FEB}      % primary blue
\definecolor{brandDark}{HTML}{0B3D91}  % deep blue
\definecolor{ink}{HTML}{1A1A1A}
\definecolor{muted}{HTML}{5B6470}
\definecolor{line}{HTML}{D9DEE6}
\definecolor{soft}{HTML}{F4F7FB}
\definecolor{ok}{HTML}{1A7F4B}
\definecolor{warn}{HTML}{B26A00}
\definecolor{danger}{HTML}{C2362B}

\graphicspath{{figures/}}
\hypersetup{
  pdftitle={Access Review — User Guide},
  pdfauthor={DashyDashboard},
  colorlinks=true, linkcolor=brand, urlcolor=brand
}

% ---------- Headings ----------
\titleformat{\section}{\Large\bfseries\color{brandDark}}{\thesection}{0.6em}{}
\titleformat{\subsection}{\large\bfseries\color{ink}}{\thesubsection}{0.6em}{}
\titlespacing*{\section}{0pt}{18pt}{8pt}
\titlespacing*{\subsection}{0pt}{12pt}{4pt}

% ---------- Header / footer ----------
\pagestyle{fancy}
\fancyhf{}
\renewcommand{\headrulewidth}{0.4pt}
\renewcommand{\footrulewidth}{0pt}
\fancyhead[L]{\small\color{muted}Access Review — User Guide}
\fancyhead[R]{\small\color{muted}\leftmark}
\fancyfoot[C]{\small\color{muted}\thepage}

% ---------- Callout boxes ----------
\newtcolorbox{tipbox}{
  enhanced, breakable, colback=soft, colframe=brand,
  boxrule=0.4pt, arc=3pt, left=10pt, right=10pt, top=7pt, bottom=7pt,
  borderline west={3pt}{0pt}{brand},
  fontupper=\small
}
\newtcolorbox{notebox}{
  enhanced, breakable, colback=white, colframe=line,
  boxrule=0.6pt, arc=3pt, left=10pt, right=10pt, top=7pt, bottom=7pt,
  fontupper=\small
}
\newtcolorbox{warnbox}{
  enhanced, breakable, colback=white, colframe=warn,
  boxrule=0.4pt, arc=3pt, left=10pt, right=10pt, top=7pt, bottom=7pt,
  borderline west={3pt}{0pt}{warn},
  fontupper=\small
}

% ---------- Screenshot helpers ----------
% \shot{file}{caption}  -> real screenshot once you add an image to ./figures
\newcommand{\shot}[2]{%
  \begin{center}
  \begin{tcolorbox}[enhanced, colback=white, colframe=line, boxrule=0.6pt,
       arc=3pt, left=4pt, right=4pt, top=4pt, bottom=4pt, width=\linewidth]
    \centering\includegraphics[width=\linewidth]{#1}\\[4pt]
    {\footnotesize\color{muted} #2}
  \end{tcolorbox}
  \end{center}}
% \shotph{caption} -> placeholder box until a screenshot is dropped in
\newcommand{\shotph}[1]{%
  \begin{center}
  \begin{tcolorbox}[enhanced, colback=soft, colframe=brand, boxrule=0pt,
       arc=4pt, left=8pt, right=8pt, top=20pt, bottom=20pt, width=\linewidth,
       borderline={0.6pt}{0pt}{brand, dashed}]
    \centering
    {\color{brand}\faImage\;\large Screenshot}\\[5pt]
    {\footnotesize\color{muted} #1}\\[3pt]
    {\scriptsize\color{muted}\itshape (drop the image into \texttt{figures/} and swap this for \texttt{{\textbackslash}shot})}
  \end{tcolorbox}
  \end{center}}

% Little inline "button" chip so steps read like the UI
\newtcbox{\btn}{on line, colback=brand, colframe=brand, coltext=white,
  boxrule=0pt, arc=3pt, left=4pt, right=4pt, top=1pt, bottom=1pt,
  fontupper=\footnotesize\bfseries}
\newtcbox{\ghost}{on line, colback=white, colframe=line, coltext=ink,
  boxrule=0.6pt, arc=3pt, left=4pt, right=4pt, top=1pt, bottom=1pt,
  fontupper=\footnotesize\bfseries}

\newcommand{\role}[1]{{\color{brandDark}\bfseries #1}}

% =====================================================================
\begin{document}

% ---------------- Title page ----------------
\begin{titlepage}
\thispagestyle{empty}
\centering
\vspace*{2.2cm}
{\color{brand}\rule{\linewidth}{2pt}}\\[6pt]
{\Huge\bfseries\color{brandDark} Access Review}\\[6pt]
{\LARGE\color{ink} User Guide}\\[10pt]
{\large\color{muted} What every person sees, and what each button does}\\[4pt]
{\color{brand}\rule{\linewidth}{2pt}}

\vfill
\begin{tcolorbox}[enhanced, width=0.85\linewidth, colback=soft, colframe=line,
   boxrule=0.6pt, arc=4pt, left=14pt, right=14pt, top=12pt, bottom=12pt]
\small
This guide walks through the dashboard from the point of view of each kind of
user --- \role{Associate}, \role{Manager}, \role{Access Management},
\role{User Directory}, and \role{Admin / Super User}. For each one it shows the
screen you can expect, the information you can see, and exactly what every button
does. No technical knowledge is needed.
\end{tcolorbox}
\vfill

{\small\color{muted} Broadridge BPO --- T+1 Access Review \quad\textbullet\quad Internal use}
\vspace*{1.5cm}
\end{titlepage}

% ---------------- Contents ----------------
\tableofcontents
\thispagestyle{fancy}
\clearpage

% =====================================================================
\section{Welcome}

The \textbf{Access Review} dashboard is where you confirm, every cycle, which
client tools you actually used. This is called \emph{attestation}: a quick,
regular sign-off that keeps everyone's access correct and audit-ready.

\subsection{What is a ``cycle''?}
A \textbf{cycle} is one review period (it runs roughly every two weeks). Each
cycle has a \textbf{due date}. During the cycle you go through your list of tools
and mark which ones you used. Managers and administrators watch how many people
have finished before the due date.

\begin{tipbox}
\faLightbulb\; \textbf{The whole job, in one sentence:} open the dashboard, mark
each tool as \emph{used} or \emph{not used} (or \emph{no access}), add a remark if
something looks wrong, then press \textbf{Submit attestation}.
\end{tipbox}

\subsection{Which role am I?}
You don't choose a role --- the system already knows who you are when you sign in.
Most people are an \role{Associate}. Some people are \emph{also} a \role{Manager},
and a small number are \role{Super Users} (Admin, GFH, GFH Delegate, or IFH). If
you have more than one role, a small set of tabs lets you switch between the views
(see Section~\ref{sec:switch}).

\begin{notebox}
\faInfoCircle\; \textbf{Roles at a glance}
\begin{itemize}[leftmargin=1.4em, itemsep=2pt]
  \item \role{Associate} --- everyone. Attest your own tools.
  \item \role{Manager} --- people who have direct reports. Track and manage your team.
  \item \role{Access Management} --- grant or remove tool access for your reports.
  \item \role{User Directory} --- Admins only. A read-only list of all users.
  \item \role{Admin Dashboard} --- Super Users (Admin / GFH / GFH Delegate / IFH).
        A department-wide overview.
\end{itemize}
\end{notebox}

% =====================================================================
\section{Signing in}

When you open the dashboard it checks your company sign-in automatically, so
usually you go straight to your dashboard --- there is nothing to type.

\shotph{The sign-in screen with the Broadridge ``B'' logo and ``Access Review'' title.}

\begin{itemize}[leftmargin=1.4em]
  \item If sign-in needs a moment, you'll see \textbf{Checking your access\ldots}
  \item If it can't confirm who you are, you'll see a short message and a
        \ghost{Try again} button. Press it to retry; if it keeps failing,
        contact the dashboard administrator.
\end{itemize}

\begin{notebox}
\faInfoCircle\; On test machines a \textbf{User name / Password} box may appear.
On the normal company site you will not see this --- your Windows sign-in is used
instead.
\end{notebox}

% =====================================================================
\section{Things you'll see everywhere}\label{sec:switch}

Most screens share the same strip across the top (the \textbf{top bar}). Once you
know it, every view feels familiar.

\shotph{The top bar: logo, cycle picker, view tabs, search, light/dark toggle, profile menu.}

\begin{longtable}{>{\raggedright}p{4.2cm} p{10.3cm}}
\toprule
\textbf{Control} & \textbf{What it does} \\
\midrule
\endhead
Broadridge logo / ``Access Review'' & Branding. Confirms you're in the right place. \\
\addlinespace
\textbf{Cycle picker} (e.g. ``\texttt{CYCLE}'') & Shows the review period you're looking at. Click it to switch to a different cycle; each entry shows how many days are left (or how long ago it ended). \\
\addlinespace
\textbf{View tabs} & Only appear if you have more than one role. Lets you jump between \ghost{Associate view}, \ghost{Manager view}, \ghost{Access}, and (Admins only) \ghost{Users}. Super Users also get \ghost{Admin Dashboard}. \\
\addlinespace
\textbf{Search box} & Filters what's on the current screen --- your tools, your team members, or departments, depending on the view. \\
\addlinespace
\textbf{Light / dark toggle} (\faMoon\,/\,\faSun) & Switches the screen between light and dark appearance. Purely visual. \\
\addlinespace
\textbf{Profile menu} & Your initials in a circle. Click for your name, role, and \ghost{Sign out}. \\
\bottomrule
\end{longtable}

\begin{tipbox}
\faLightbulb\; \textbf{Switching roles:} if you're a Manager or Super User, use the
view tabs in the top bar to move between your Associate, Manager, Access, and Admin
screens. Switching never loses your work --- everything saves as you go.
\end{tipbox}

% =====================================================================
\section{Role 1 --- Associate (everyone)}

This is the main screen and the one everyone uses. It lists the clients you work
with and, under each, the tools you have access to. Your job is to say whether you
used each tool this cycle.

\shotph{The Associate dashboard: greeting, progress card, and a client with its tools expanded.}

\subsection{The top of the screen}
\begin{itemize}[leftmargin=1.4em]
  \item A friendly greeting and a line such as \emph{``You have 4 tools left to
        attest before <due date>.''}
  \item Four quick numbers: \textbf{Attested} (done), \textbf{Remaining} (still to
        do), \textbf{Clients}, and \textbf{Days left}.
  \item A \textbf{Cycle progress} card showing how many of your tools are done out
        of the total, a percentage, a progress bar, and the
        \btn{Submit attestation} button.
\end{itemize}

\subsection{Your clients and tools}
Each client is a coloured card. Click a card's header (or the chevron \faChevronDown
on the right) to open or close it. The header shows the tool count, how many you've
marked as used, and how many are still pending.

When open, you see a table with four columns:

\begin{longtable}{>{\raggedright}p{3.2cm} p{11.3cm}}
\toprule
\textbf{Column} & \textbf{What to do with it} \\
\midrule
\endhead
\textbf{Tool} & The name of the tool. Nothing to click. \\
\addlinespace
\textbf{Had access?} & A \ghost{Yes}\,/\,\ghost{No} switch. Leave it on \textbf{Yes} if you really had access to this tool. Choose \textbf{No} if you should \emph{not} have had access --- this flags a possible access dispute for your manager. \\
\addlinespace
\textbf{Did you use?} & A \ghost{Yes}\,/\,\ghost{No} switch --- did you actually use this tool this cycle? This is greyed out if you set \emph{Had access?} to \textbf{No} (you can't ``use'' a tool you shouldn't have). \\
\addlinespace
\textbf{Remark} & A small \ghost{\faCommentDots\ Add remark} button. Use it to leave a note --- for example, why a tool was or wasn't used. If a note already exists, its text shows on the button. \\
\bottomrule
\end{longtable}

\begin{notebox}
\faInfoCircle\; \textbf{Everything saves instantly.} Each time you flip a switch it
is saved straight away --- there's no separate ``save'' step for individual tools.
A tool counts as done once you've answered \emph{Did you use?} \textbf{or} marked
\emph{Had access? = No}.
\end{notebox}

\subsection{Adding a remark}
Clicking \ghost{\faCommentDots\ Add remark} opens a small window for that one tool.

\shotph{The remark window with a text box, character counter, Cancel and Save remark buttons.}

\begin{itemize}[leftmargin=1.4em]
  \item Type your note (up to 500 characters).
  \item \btn{Save remark} stores it; \ghost{Cancel} closes without changes.
  \item Leaving the box empty and saving \textbf{clears} the remark.
  \item Remarks are visible to your manager and to audit.
\end{itemize}

\subsection{Submitting}
When you've gone through your tools, press \btn{Submit attestation} in the progress
card. You'll get a short confirmation message. You can keep updating and submit
again if something changes before the due date.

\begin{warnbox}
\faExclamationTriangle\; If you press \textbf{Submit} without marking anything, the
dashboard reminds you to mark at least one tool first.
\end{warnbox}

\begin{tipbox}
\faLightbulb\; \textbf{Search tip:} use the top-bar search to jump to a tool by
name --- handy if you have many clients.
\end{tipbox}

% =====================================================================
\section{Role 2 --- Manager}

If you have direct reports, the \ghost{Manager view} tab appears in the top bar.
This view is about \emph{watching your team's progress} --- you don't attest on
their behalf.

\shotph{Manager view: summary cards, the team list on the left, and a selected member on the right.}

\subsection{What you see}
\begin{itemize}[leftmargin=1.4em]
  \item A short \textbf{Manager view} banner naming the current cycle.
  \item Four summary cards: \textbf{Direct reports}, \textbf{Submitted},
        \textbf{In progress}, \textbf{Not started}.
  \item \textbf{Team completion} (left) --- a list of your people with a progress
        bar and a status badge each.
  \item \textbf{Selected member} (right) --- details for whoever you click.
\end{itemize}

\subsection{The team list (left panel)}
\begin{longtable}{>{\raggedright}p{4.0cm} p{10.5cm}}
\toprule
\textbf{Control} & \textbf{What it does} \\
\midrule
\endhead
Filter pills & \ghost{Everyone}, \ghost{Not started}, \ghost{In progress}, \ghost{Submitted} --- show only people in that state. \\
\addlinespace
A team-member row & Shows their photo initials, name, ID, progress (e.g. ``6/8''), and a status badge. Click a row to load their detail on the right. \\
\addlinespace
Top-bar search & Filter the list by name or associate ID. \\
\bottomrule
\end{longtable}

\subsection{Member detail (right panel)}
For the selected person you'll see their completion percentage, status, and a
\textbf{per-client progress} breakdown. If they reported any access problems, an
\textbf{Access disputes} list appears in red with the tool, client, and their remark.

\subsection{Exporting disputes}
If anyone on your team flagged an access dispute (a \emph{Had access? = No}), a red
banner appears at the top with an \btn{\faDownload\ Export disputes} button. Click
it to download a file of those records to follow up on.

\shotph{The red ``access disputes this cycle'' banner with the Export disputes button.}

% =====================================================================
\section{Role 3 --- Access Management}

The \ghost{Access} tab is where managers \textbf{grant, time-box, or remove} tool
access for their direct reports.

\shotph{Access management: team list on the left, grant form and active-access table on the right.}

\subsection{Choosing a person}
The left sidebar lists your direct reports with a search box at the top. Click a
name to manage that person's access on the right.

\subsection{Granting new access}
The \textbf{Grant new access} form has four fields:

\begin{longtable}{>{\raggedright}p{3.6cm} p{10.9cm}}
\toprule
\textbf{Field} & \textbf{What to enter} \\
\midrule
\endhead
\textbf{Client} & Pick the client from the drop-down. \\
\addlinespace
\textbf{Tool} & Pick the tool. (Available only after a client is chosen.) \\
\addlinespace
\textbf{Access from} & The start date. Defaults to today. \\
\addlinespace
\textbf{Access to (optional)} & An end date. Leave blank for open-ended access. \\
\bottomrule
\end{longtable}

Press \btn{\faPlus\ Grant access} to save. The new access appears in the table below.

\subsection{Managing existing access}
The \textbf{Active access} table lists everything the person currently has:

\begin{longtable}{>{\raggedright}p{3.0cm} p{11.5cm}}
\toprule
\textbf{Column} & \textbf{What it does} \\
\midrule
\endhead
\textbf{Client} / \textbf{Tool} & Which access this row is for. \\
\addlinespace
\textbf{From} & The date access started. \\
\addlinespace
\textbf{To} & An editable end date --- change it to time-box the access. \\
\addlinespace
\textbf{Actions} & \ghost{Clear} removes the end date (makes it open-ended); \ghost{Revoke now} ends the access today. \\
\bottomrule
\end{longtable}

\begin{warnbox}
\faExclamationTriangle\; \textbf{Revoke now} takes effect immediately. Use it when
someone should lose access to a tool right away.
\end{warnbox}

% =====================================================================
\section{Role 4 --- User Directory (Admins only)}

Admins get a \ghost{Users} tab: a \textbf{read-only} directory of everyone in the
system. There are no edit buttons here --- it's for looking people up.

\shotph{The User Directory: a searchable table of ID, Full Name, User Name, Department, Manager.}

\begin{itemize}[leftmargin=1.4em]
  \item A \textbf{search box} filters by name, ID, user name, department, or manager.
  \item The table columns are \textbf{ID}, \textbf{Full Name}, \textbf{User Name},
        \textbf{Department}, and \textbf{Manager}.
  \item Clicking a row simply highlights it.
\end{itemize}

% =====================================================================
\section{Role 5 --- Admin Dashboard (Super Users)}

Super Users (roles \role{Admin}, \role{GFH}, \role{GFH Delegate}, \role{IFH}) get
the \ghost{Admin Dashboard} --- a department-wide overview of how attestation is
going.

\shotph{Admin Dashboard: sidebar, KPI band across the top, and a grid of department cards with donut charts.}

\subsection{What's on screen}
\begin{itemize}[leftmargin=1.4em]
  \item A \textbf{sidebar} on the left with the \ghost{Overview} item, your
        ``Logged in as'' badge, \ghost{Sign Out}, and (if you're also a manager) a
        \ghost{Team View} shortcut back to the manager screen.
  \item A \textbf{top header} with the page title, a \textbf{Search departments}
        box, the \textbf{Cycle picker}, a \textbf{notifications} bell, the
        light/dark toggle, and your profile menu.
  \item A \textbf{KPI band} of five numbers: \textbf{Associates},
        \textbf{Associate-Tools}, \textbf{Submitted} (with \%),
        \textbf{Pending}, and \textbf{At Risk Depts}.
  \item A grid of \textbf{department cards}, each with a donut chart, the
        department's GFH, a status badge, and a submitted-of-total count.
\end{itemize}

\subsection{Department cards}
\begin{itemize}[leftmargin=1.4em]
  \item The \faInfoCircle\ icon next to ``GFH'' opens a small card with the Global
        Functional Head's name, email, and office.
  \item Click anywhere else on a card to \textbf{drill into} that department.
  \item Some cards show a per-client breakdown with mini progress bars underneath.
\end{itemize}

\subsection{Drilling into a department}
\shotph{Department drill-down: back button, client filter pills, the department header, the Exception Exports card, and manager cards.}

\begin{longtable}{>{\raggedright}p{4.2cm} p{10.3cm}}
\toprule
\textbf{Control} & \textbf{What it does} \\
\midrule
\endhead
\ghost{\faArrowLeft\ All Departments} & Goes back to the department grid. \\
\addlinespace
\textbf{Filter by client} pills & Narrow the figures to a single client, or \ghost{All clients}. \\
\addlinespace
\textbf{Completion breakdown} bar & Shows Submitted / In Progress / Not Started for the department. \\
\addlinespace
\textbf{Exception Exports} & Two download buttons: \ghost{Export} next to \emph{Incomplete submissions} (people who haven't finished) and \ghost{Export} next to \emph{Access disputes} (flagged access). \\
\addlinespace
\textbf{Manager cards} & One per manager in the department, each with a donut and a quick info popover (\faInfoCircle) for their email. \\
\bottomrule
\end{longtable}

\subsection{Adding a tool}
Admins, GFH, and GFH Delegates see an \btn{\faPlus\ Add Tool} button in the header.

\shotph{The Add Tool window with Client, Department, and Tool Name fields.}

\begin{itemize}[leftmargin=1.4em]
  \item Choose the \textbf{Client}, choose the \textbf{Department}, and type the
        \textbf{Tool Name}.
  \item \btn{Add Tool} creates it; \ghost{Cancel} closes without saving.
  \item Once added, the tool is available for that department's managers to assign.
\end{itemize}

\subsection{Notifications}
The bell (\faBell) in the header lists departments that are behind as the due date
approaches, plus a cycle reminder. A small number badge shows how many need
attention; \ghost{Mark all read} clears it.

\subsection{How the view differs by role}
The Admin Dashboard adjusts to who you are:

\begin{longtable}{>{\raggedright}p{3.2cm} p{4.8cm} p{5.7cm}}
\toprule
\textbf{Role} & \textbf{Which departments} & \textbf{Extras} \\
\midrule
\endhead
\role{Admin} & All departments & Full overview, \ghost{Add Tool}, \ghost{Users} directory tab \\
\addlinespace
\role{GFH Delegate} & All departments & Same wide view as Admin, plus \ghost{Add Tool} \\
\addlinespace
\role{GFH} & Only your own department(s) & \ghost{Add Tool}; title shows your name and role \\
\addlinespace
\role{IFH} & Only your own department(s) & Read/track focus; no \ghost{Add Tool} \\
\bottomrule
\end{longtable}

% =====================================================================
\section{Quick reference}

\subsection{Who can see what}
\begin{longtable}{>{\raggedright}p{4.2cm} c c c c c}
\toprule
\textbf{View / Action} & \textbf{Assoc.} & \textbf{Mgr} & \textbf{Access} & \textbf{Admin} & \textbf{GFH/IFH} \\
\midrule
\endhead
Attest my own tools        & \faCheck & \faCheck & \faCheck & \faCheck & \faCheck \\
Add remarks                & \faCheck & \faCheck & \faCheck & \faCheck & \faCheck \\
See my team's progress     &          & \faCheck &          & \faCheck & \faCheck \\
Export team disputes       &          & \faCheck &          & \faCheck & \faCheck \\
Grant / revoke access      &          &          & \faCheck &          &          \\
User Directory             &          &          &          & \faCheck &          \\
Department overview        &          &          &          & \faCheck & \faCheck \\
Add a tool                 &          &          &          & \faCheck & GFH only \\
\bottomrule
\end{longtable}

\begin{notebox}
\faInfoCircle\; The tabs you see depend on your roles. Most people only ever see
the Associate view. Managers also see Manager and Access. Super Users additionally
see the Admin Dashboard, and Admins also see the Users directory.
\end{notebox}

\subsection{Glossary}
\begin{longtable}{>{\raggedright}p{3.6cm} p{10.9cm}}
\toprule
\textbf{Term} & \textbf{Meaning} \\
\midrule
\endhead
\textbf{Attestation} & Your sign-off that says which tools you used this cycle. \\
\textbf{Cycle} & One review period (about two weeks) with its own due date. \\
\textbf{Had access?} & Whether you legitimately had access to a tool. ``No'' raises an access dispute. \\
\textbf{Did you use?} & Whether you actually used the tool during the cycle. \\
\textbf{Access dispute} & A flag that someone had access they shouldn't have --- surfaced to managers and admins. \\
\textbf{Remark} & A free-text note attached to one tool, visible to your manager and audit. \\
\textbf{GFH} & Global Functional Head --- senior owner of a department. \\
\textbf{IFH} & Functional head with a department-scoped view. \\
\textbf{Super User} & An Admin, GFH, GFH Delegate, or IFH with the Admin Dashboard. \\
\bottomrule
\end{longtable}

\vspace{8pt}
\begin{tipbox}
\faLifeRing\; \textbf{Need help?} If something doesn't look right --- a missing tool,
access you don't recognise, or a screen that won't load --- use the
\emph{Had access? = No} flag and a remark where relevant, and contact your manager
or the dashboard administrator.
\end{tipbox}

\end{document}
